本研究で提案した手法が、実際のソフトウェア開発現場でどのように貢献し、開発者の意思決定を支援するかについて、具体的なユースケースを用いて説明する。

\paragraph{迅速なリリースが求められる大規模Webサービスの開発}
ある大規模Webサービスを開発しているチームを想定する。このチームではCI/CDパイプラインを導入しており、1日に数十件のコミットが発生している。チームが抱える人員に対してシステムの規模が大きいため、すべての変更を詳細にレビューすることは現実的ではなく、多くの場合、変更行数の少ない小規模な変更のレビューが簡略化されがちである。ここで本手法を適用すると、以下のようなプロセスで品質向上が図られる。

初めに、開発者が「特定のメソッド内の数トークンのみを変更した小規模なコミット」をプッシュした際、本モデルがそれを「高リスク」と判定し、レビューツール上に警告を表示する。従来の変更規模に基づく評価では見逃されていた、局所的な変更が全体に及ぼすリスクを数値として提示する。

続いて、プロジェクトリーダーは、本手法が出力する「密度 $D_{i}$（単位労力あたりの欠陥発見期待値）」に基づき、その日のレビュー対象を優先順位付けする。実験結果（5.7節）で示した通り、総労力のわずか20\%を「密度の高いコミット」に集中させるだけで、潜在的な欠陥の約7割を早期に特定できる。

また、特徴量重要度の分析により、なぜその変更が「高リスク」と判定されたのか（例: 特定の複雑なメソッドのトークン変化が激しいため）を開発者にフィードバックできる。これにより、開発者は自身の修正が既存コードの整合性にどう影響したかを再認識し、修正の質を自発的に向上させることが可能となる。

本手法の最大の長所は、人間の直感に依存しないレビュー労力の分配を可能にする点にある。従来、「小規模な変更は安全である」と過信されていた領域に潜むリスクを、時系列変化と複数の構成要素からなるメトリクスを組み合わせることで明らかにした。これにより、開発チームは、どこを重点的に見るべきか分からないという悩みから解放され、新機能の開発やリファクタリングに注力できるようになる。